\documentclass[aspectratio=169]{beamer}
\usetheme{metropolis}
\usepackage{appendixnumberbeamer}
\usepackage{booktabs, hyperref}

\input{mgmt675-style}

\usepackage{tikz}
\usetikzlibrary{shapes.geometric, arrows.meta, positioning, calc}

% Title info
\subtitle{MGMT 675: Generative AI for Finance}
\title{Module 3: Mean-Variance Analysis}
\author{Kerry Back}
\date{}

\begin{document}

\maketitle

\begin{frame}{Claude.ai Is an Agent}
Claude.ai is an agent connected to a \textbf{Python environment} --- it writes and runs code in a cloud sandbox to answer your questions.
\vspace{0.3cm}

Claude.ai also has access to other tools:
\begin{itemize}
  \item \textbf{Web browser} --- search the web and read pages
  \item \textbf{File system} --- read and write files you upload or create
  \item \textbf{Artifacts} --- build interactive apps and charts in a side panel
  \item \textbf{Document generation} --- produce Excel, Word, and PowerPoint files
\end{itemize}
\vspace{0.3cm}
When you ask Claude.ai to analyze data, it plans the code, executes it in the Python environment, observes the result, and iterates --- the same agent loop from Module~1.
\end{frame}

% ============================================================
% PORTFOLIO OPTIMIZATION
% ============================================================
\section{Portfolio Optimization}

\begin{frame}{The Problem}
Given expected returns, risks (standard deviations), and correlations for multiple assets, plus a risk-free rate:

\begin{itemize}
  \item What is the \textbf{best portfolio}?
  \item What does the set of all possible portfolios look like?
  \item How do real-world constraints change the answer?
\end{itemize}

\vspace{0.3cm}
Mean-variance analysis provides the framework.  AI handles the computation.
\end{frame}

\begin{frame}{Key Concepts}
\begin{columns}[T]
\begin{column}{0.48\textwidth}
\begin{shadedbox}[title=\textbf{Portfolios}]
\begin{itemize}\small
  \item \textbf{Tangency portfolio:} the risky portfolio with the highest Sharpe ratio (excess return per unit of risk)
  \item \textbf{Global minimum variance (GMV):} the risky portfolio with the lowest possible risk
  \item \textbf{Efficient frontier:} the set of risky portfolios with the highest return for each level of risk
\end{itemize}
\end{shadedbox}
\end{column}
\begin{column}{0.48\textwidth}
\begin{shadedbox}[title=\textbf{Capital Allocation Line}]
\begin{itemize}\small
  \item The straight line from the risk-free rate through the tangency portfolio
  \item Represents all combinations of the risk-free asset and the tangency portfolio
  \item The best attainable risk-return tradeoff
\end{itemize}
\end{shadedbox}
\end{column}
\end{columns}
\end{frame}

\begin{frame}{Adding Real-World Constraints}
\begin{columns}[T]
\begin{column}{0.48\textwidth}
\textbf{Common constraints}
\begin{itemize}
  \item No short sales (all weights $\geq 0$)
  \item Maximum position (e.g., $\leq 40\%$ per asset)
  \item Minimum position (e.g., $\geq 2\%$ if held)
  \item Margin requirement (sum of absolute values $\leq 2$)
\end{itemize}
\end{column}
\begin{column}{0.48\textwidth}
\textbf{What happens}
\begin{itemize}
  \item The efficient frontier shifts \alert{inward} (lower returns and/or higher risk)
  \item The tangency portfolio changes
  \item The Sharpe ratio cannot improve with constraints---it can only stay the same or get worse
\end{itemize}
\end{column}
\end{columns}
\vspace{0.3cm}
\alert{The constrained frontier is still useful: it reflects what you can actually implement.}
\end{frame}

\begin{frame}{Solution Methods}
\begin{itemize}
  \item \textbf{Solver (numerical optimization)}
    \begin{itemize}
      \item Maximize Sharpe ratio (tangency portfolio)
      \item Minimize risk for a target return (efficient frontier)
      \item Minimize risk (global minimum variance portfolio)
    \end{itemize}
  \item \textbf{Analytic/algebraic solution}
    \begin{itemize}
      \item Solve systems of linear equations for tangency, GMV, and frontier
      \item Available only when there are no constraints
    \end{itemize}
\end{itemize}

\vspace{0.3cm}
Python solver options: \texttt{scipy.minimize}, \texttt{cvxopt}, \texttt{cvxpy}---AI chooses and writes the code.

\vfill
\alert{Ask Claude:} \small Discuss the advantages and disadvantages of these solver options for mean-variance analysis with inequality constraints.
\end{frame}

\begin{frame}{In-Class Exercise: Data}
We will use monthly returns for five ETFs:
\vspace{0.2cm}
\begin{itemize}
  \item \textbf{SPY} --- S\&P 500
  \item \textbf{GLD} --- Gold
  \item \textbf{UUP} --- U.S.\ Dollar Index
  \item \textbf{IEF} --- 7--10 Year Treasuries
  \item \textbf{EFA} --- International Developed Markets
\end{itemize}
\vspace{0.2cm}
Data: April 2007 -- present (228 months).  Risk-free rate: 3.75\% (current 1-month T-bill).
\vspace{0.3cm}

Ask Claude to use \texttt{yfinance} to download monthly prices for these five tickers, compute monthly returns, and set $r_f = 3.75\%$ annualized.
\end{frame}

\begin{frame}{In-Class Exercise: Tangency Portfolio (Unrestricted)}
\begin{enumerate}
  \item Ask Claude to compute the \textbf{tangency portfolio} allowing short sales.  Report the weights, expected return, standard deviation, and Sharpe ratio.
  \item Ask Claude: \textit{``Show me the code you used and explain it step by step.''}
\end{enumerate}
\vspace{0.3cm}
\alert{Understanding the code matters.}  You should be able to verify that the optimization is set up correctly --- what is the objective function?  What are the constraints?
\end{frame}

\begin{frame}{In-Class Exercise: No Short Sales}
\begin{enumerate}
  \item Ask Claude to compute the \textbf{tangency portfolio with no short sales} (all weights $\geq 0$).  Report the weights, expected return, standard deviation, and Sharpe ratio.
  \item Again ask: \textit{``Show me the code you used and explain it.''}
  \item Compare with the unrestricted result --- how did the weights and Sharpe ratio change?
\end{enumerate}
\end{frame}

\begin{frame}{In-Class Exercise: Frontier and CAL}
Assuming \textbf{no short sales}:
\begin{enumerate}
  \item Ask Claude to compute the \textbf{tangency portfolio} and the \textbf{efficient frontier} from the five risky assets.
  \item Ask Claude to \textbf{plot} the efficient frontier and the \textbf{capital allocation line} (the line from $r_f$ through the tangency portfolio).
\end{enumerate}
\end{frame}

\end{document}
